\documentclass{article}

% ---------------------------------------------------------------------------
% Formato NeurIPS. Necesitas el archivo de estilo neurips_2024.sty junto a
% este main.tex (ver COMO_COMPILAR.md). En Overleaf, parte de la plantilla
% oficial "NeurIPS 2024" y pega este contenido.
%
% Si NO tienes el .sty, comenta la línea \usepackage[final]{neurips_2024}
% y descomenta el bloque "FALLBACK" más abajo para compilar con article.
% ---------------------------------------------------------------------------
\usepackage[final]{neurips_2024}

% --- FALLBACK (sin neurips_2024.sty): descomenta estas 3 líneas ---
% \usepackage[margin=1in]{geometry}
% \usepackage{times}
% \renewcommand{\headeright}{}
% ------------------------------------------------------------------

\usepackage[utf8]{inputenc}
\usepackage[T1]{fontenc}
\usepackage[spanish,es-noindentfirst]{babel}
\usepackage{hyperref}
\usepackage{url}
\usepackage{booktabs}
\usepackage{amsmath}
\usepackage{amsfonts}
\usepackage{nicefrac}
\usepackage{microtype}
\usepackage{xcolor}
\usepackage{graphicx}
\usepackage{subcaption}

\graphicspath{{figures/}}

% Comando para marcar texto pendiente de completar
\newcommand{\TODO}[1]{\textcolor{red}{[\textbf{COMPLETAR}: #1]}}

\title{Clasificación de Llantas Dañadas mediante Reconocimiento de Textura:\\
       CNN desde Cero frente a \emph{Transfer Learning}, con Interpretabilidad
       Grad-CAM y Manejo del Desbalance de Clases}

% Reemplaza con los integrantes del grupo.
\author{%
  Nombre Apellido \\
  Afiliación / Curso \\
  \texttt{correo@dominio} \\
  \And
  Nombre Apellido \\
  Afiliación / Curso \\
  \texttt{correo@dominio} \\
}

\begin{document}

\maketitle

\begin{abstract}
La inspección manual del estado de las llantas de vehículos es costosa, subjetiva
y difícil de escalar. Proponemos un sistema de clasificación binaria basado en redes
neuronales convolucionales (CNN) que distingue entre llantas en buen estado y
llantas dañadas a partir de la textura superficial. Comparamos de forma rigurosa una
CNN entrenada desde cero contra modelos de \emph{transfer learning}
(ResNet-50 y EfficientNet-B0 preentrenados en ImageNet), evaluando con métricas
sensibles al desbalance (precisión, \emph{recall}, F1 y AUC-ROC). Analizamos el efecto
del desbalance de clases y cuantificamos el impacto de tres estrategias de mitigación
(pesos de clase, sobremuestreo y pérdida focal). Mediante Grad-CAM interpretamos qué
patrones de textura el modelo asocia al daño, y discutimos cualitativamente los modos
de fallo. \TODO{resumir aquí el resultado principal en 1--2 frases con cifras, p.\,ej.
``el mejor modelo alcanza F1$=$0.XX y AUC$=$0.XX''.}
\end{abstract}

% ===========================================================================
\section{Introducción y Motivación}
% ===========================================================================
El desgaste y el daño en las llantas constituyen un riesgo crítico para la seguridad
vial. La inspección visual humana no escala a entornos industriales y presenta
variabilidad entre inspectores. Un clasificador automático fiable podría servir de
apoyo al tamizaje, priorizando unidades para revisión experta.

Este trabajo aborda la detección de daño como un problema de clasificación binaria de
imágenes de textura. Nuestras contribuciones son: (i) la implementación de una CNN
propia desde cero y su comparación rigurosa contra \emph{backbones} preentrenados;
(ii) un estudio del desbalance de clases con varias estrategias de mitigación;
(iii) un análisis de interpretabilidad con Grad-CAM que vincula las decisiones del
modelo con regiones de textura; y (iv) un análisis cualitativo de fallos.
\TODO{ajustar las contribuciones a lo que efectivamente obtuviste.}

% ===========================================================================
\section{Trabajos Relacionados}
% ===========================================================================
Las CNN profundas transformaron la clasificación de imágenes a partir de AlexNet
\cite{krizhevsky2012alexnet} sobre ImageNet \cite{deng2009imagenet}. La normalización
por lotes \cite{ioffe2015batchnorm} estabilizó el entrenamiento de redes profundas, y
las conexiones residuales de ResNet \cite{he2016resnet} permitieron entrenar
arquitecturas muy profundas. EfficientNet \cite{tan2019efficientnet} propuso un escalado
compuesto que mejora la relación precisión/cómputo. Para problemas desbalanceados, la
pérdida focal \cite{lin2017focal} reduce la contribución de los ejemplos fáciles y
enfatiza los difíciles. En interpretabilidad, Grad-CAM \cite{selvaraju2017gradcam}
produce mapas de localización a partir de los gradientes de la última capa
convolucional. \TODO{añade al menos una referencia específica de inspección de
defectos/textura o de detección de daño en llantas, y la cita del propio dataset
(mínimo 5 referencias en total).}

% ===========================================================================
\section{Metodología}
% ===========================================================================
\subsection{Datos y particiones}
Empleamos el conjunto \emph{Tire Texture Image Recognition} (Kaggle), con imágenes
organizadas en carpetas por clase, mapeadas a las etiquetas binarias
\texttt{good} (0) y \texttt{damaged} (1). Realizamos una partición estratificada
70\%/15\%/15\% en entrenamiento/validación/prueba con semilla fija para
reproducibilidad. \TODO{reporta el número de imágenes por clase y por partición
(tomados de \texttt{outputs/<run>/data\_info.json}).}

\subsection{Arquitecturas}
Comparamos tres modelos. (1)~Una \textbf{CNN desde cero} de cuatro bloques
convolucionales (Conv--BN--ReLU $\times 2$ + \emph{max-pooling}), \emph{global average
pooling} y una cabeza totalmente conectada con \emph{dropout}. (2)~\textbf{ResNet-50}
\cite{he2016resnet} preentrenada en ImageNet con la capa final reemplazada por un único
logit. (3)~\textbf{EfficientNet-B0} \cite{tan2019efficientnet} con un reemplazo análogo
de la cabeza. Todos los modelos emiten un solo logit (clase positiva $=$ daño).

\subsection{Funciones de pérdida y desbalance de clases}
Como pérdida base usamos entropía cruzada binaria con logits (BCE). Para el desbalance
evaluamos: (a) \textbf{pesos de clase} vía \texttt{pos\_weight} $= N_{\text{neg}}/N_{\text{pos}}$;
(b) \textbf{sobremuestreo} con \texttt{WeightedRandomSampler}; y (c) \textbf{pérdida focal}
\cite{lin2017focal},
\begin{equation}
\mathcal{L}_{\text{focal}} = -\alpha_t\,(1-p_t)^{\gamma}\,\log(p_t),
\end{equation}
con $\alpha=0.25$ y $\gamma=2.0$ por defecto.

\subsection{Protocolo de entrenamiento}
Optimizador AdamW, planificador de tasa de aprendizaje por \emph{cosine annealing},
entrenamiento en precisión mixta (AMP) y \emph{early stopping} según el F1 de
validación. \TODO{especifica épocas, tamaño de lote, tasa de aprendizaje y aumento de
datos finales (de \texttt{config.json}).} El aumento de datos contempla recorte
aleatorio, volteos, rotaciones y \emph{jitter} de color; la variante ``strong'' añade
desenfoque y \emph{random erasing}.

% ===========================================================================
\section{Experimentos}
% ===========================================================================
Diseñamos tres bloques experimentales: (A) comparación de arquitecturas;
(B) ablación sobre tres factores (función de pérdida, intensidad del aumento de datos
y profundidad del \emph{backbone}); y (C) estudio del desbalance de clases. Todas las
métricas se reportan sobre el conjunto de prueba retenido.

% ===========================================================================
\section{Resultados}
% ===========================================================================
\subsection{Comparación de arquitecturas}
La Tabla~\ref{tab:arch} resume el desempeño. \TODO{completa con
\texttt{outputs/comparison\_architectures.csv} y comenta qué arquitectura gana y por
qué (capacidad vs.\ datos disponibles, beneficio del preentrenamiento).}

\begin{table}[ht]
  \centering
  \caption{Comparación de arquitecturas en el conjunto de prueba.}
  \label{tab:arch}
  \begin{tabular}{lccccc}
    \toprule
    Modelo & Exactitud & Precisión & Recall & F1 & AUC-ROC \\
    \midrule
    CNN desde cero       & -- & -- & -- & -- & -- \\
    ResNet-50 (transfer) & -- & -- & -- & -- & -- \\
    EfficientNet-B0      & -- & -- & -- & -- & -- \\
    \bottomrule
  \end{tabular}
\end{table}

La Figura~\ref{fig:cm_roc} muestra la matriz de confusión y la curva ROC del mejor
modelo. \TODO{interpreta los falsos negativos (daño no detectado) frente a los falsos
positivos.}

\begin{figure}[ht]
  \centering
  \begin{subfigure}{0.42\textwidth}
    \includegraphics[width=\linewidth]{confusion_matrix.png}
    \caption{Matriz de confusión.}
  \end{subfigure}\hfill
  \begin{subfigure}{0.42\textwidth}
    \includegraphics[width=\linewidth]{roc_curve.png}
    \caption{Curva ROC.}
  \end{subfigure}
  \caption{Desempeño del mejor modelo en prueba. \TODO{copiar las imágenes desde
  \texttt{outputs/<run>/} a la carpeta \texttt{figures/}.}}
  \label{fig:cm_roc}
\end{figure}

\subsection{Estudio de ablación}
La Tabla~\ref{tab:ablation} reporta el efecto de cada factor.
\TODO{completa con \texttt{outputs/ablation.csv}.}

\begin{table}[ht]
  \centering
  \caption{Ablación: pérdida, aumento de datos y profundidad del backbone.}
  \label{tab:ablation}
  \begin{tabular}{llccc}
    \toprule
    Factor & Variante & Recall & F1 & AUC-ROC \\
    \midrule
    Pérdida   & BCE            & -- & -- & -- \\
    Pérdida   & Focal          & -- & -- & -- \\
    Aumento   & none           & -- & -- & -- \\
    Aumento   & standard       & -- & -- & -- \\
    Aumento   & strong         & -- & -- & -- \\
    Backbone  & ResNet-18      & -- & -- & -- \\
    Backbone  & ResNet-50      & -- & -- & -- \\
    \bottomrule
  \end{tabular}
\end{table}

\subsection{Desbalance de clases}
\TODO{indica el ratio de desbalance (de \texttt{data\_info.json}).} La
Tabla~\ref{tab:imbalance} compara las estrategias de mitigación frente al modelo sin
mitigación. \TODO{completa con \texttt{outputs/imbalance\_study.csv} y discute el
\emph{trade-off}: las estrategias que suben el \emph{recall} de la clase de daño
suelen costar precisión.}

\begin{table}[ht]
  \centering
  \caption{Efecto de las estrategias de mitigación del desbalance.}
  \label{tab:imbalance}
  \begin{tabular}{lcccc}
    \toprule
    Estrategia & Precisión & Recall & F1 & AUC-ROC \\
    \midrule
    Sin mitigación & -- & -- & -- & -- \\
    Pesos de clase & -- & -- & -- & -- \\
    Sobremuestreo  & -- & -- & -- & -- \\
    Pérdida focal  & -- & -- & -- & -- \\
    \bottomrule
  \end{tabular}
\end{table}

\subsection{Interpretabilidad (Grad-CAM)}
La Figura~\ref{fig:gradcam} muestra mapas de activación para muestras de ambas clases.
\TODO{describe qué patrones de textura (grietas, deformaciones, bordes irregulares)
activan la predicción de daño, y si las activaciones son plausibles.}

\begin{figure}[ht]
  \centering
  \includegraphics[width=0.7\textwidth]{gradcam.png}
  \caption{Mapas Grad-CAM sobre ejemplos \texttt{good} y \texttt{damaged}.}
  \label{fig:gradcam}
\end{figure}

\subsection{Análisis cualitativo de fallos}
La Figura~\ref{fig:failures} recopila los ejemplos peor clasificados.
\TODO{comenta al menos cinco casos con hipótesis de causa (iluminación, oclusión,
ambigüedad de la etiqueta, daño sutil, textura atípica).}

\begin{figure}[ht]
  \centering
  \includegraphics[width=0.85\textwidth]{failures.png}
  \caption{Ejemplos mal clasificados (falsos negativos y falsos positivos).}
  \label{fig:failures}
\end{figure}

% ===========================================================================
\section{Discusión y Limitaciones}
% ===========================================================================
\TODO{discute de forma crítica:} tamaño y representatividad del dataset; posible sesgo
en las condiciones de captura; brecha entre el desempeño en prueba y el esperado en
producción; riesgos de despliegue (un falso negativo deja pasar una llanta peligrosa);
y dependencia de los pesos preentrenados de ImageNet, cuyo dominio difiere del de
texturas de llantas.

% ===========================================================================
\section{Conclusión}
% ===========================================================================
\TODO{resume el hallazgo principal, la mejor configuración y una línea de trabajo
futuro (p.\,ej.\ segmentación del daño o detección multiclase de tipos de defecto).}

\bibliographystyle{plain}
\bibliography{references}

\end{document}
